\documentclass[12pt,a4paper]{article}

\usepackage[T1]{fontenc}
\usepackage[utf8]{inputenc}
\usepackage[spanish]{babel}
\usepackage{times}
\usepackage{setspace}
\usepackage{geometry}
\usepackage{graphicx}
\usepackage{booktabs}
\usepackage{longtable}
\usepackage{array}
\usepackage{hyperref}
\usepackage{enumitem}
\usepackage{float}
\usepackage{xcolor}
\usepackage{caption}
\usepackage{subcaption}
\usepackage{tikz}
\usetikzlibrary{positioning,arrows.meta,shapes.geometric,fit}

\geometry{margin=2.54cm}
\hypersetup{hidelinks}
\setstretch{1.75}
\setlength{\parindent}{0.5in}
\setlength{\parskip}{0pt}
\captionsetup{font=small,labelfont=bf}
\addto\captionsspanish{\renewcommand{\figurename}{Figure}}

\newcommand{\EvidenceImage}[1]{%
    \IfFileExists{#1}{%
        \includegraphics[width=0.98\linewidth,height=0.16\textheight,keepaspectratio]{#1}
    }{%
        \fbox{\parbox{0.92\linewidth}{\centering\textbf{Imagen no encontrada}\\\texttt{#1}}}
    }%
}

\title{Aplicacion de Control Academico\\Informe Tecnico (Formato Academico APA Adaptado)}
\author{%
Jesus Alvarez Sombrerero (177516)\\
Camila Garcia Alvarez (178004)\\
Celeste Isabel Alonso Garcia (173291)
}
\date{\today}

\begin{document}

\begin{titlepage}
    \centering
    \vspace*{2.5cm}
    {\Large Universidad de las Americas Puebla\\[0.4cm]}
    {\large Ingenieria en Sistemas Computacionales\\[1.8cm]}
    {\LARGE \textbf{Aplicacion de Control Academico}}\\[0.3cm]
    {\Large Informe tecnico con enfoque de arquitectura MVVM}\\[1.6cm]

    \begin{flushleft}
    \textbf{Integrantes:}\\
    Jesus Alvarez Sombrerero (177516)\\
    Camila Garcia Alvarez (178004)\\
    Celeste Isabel Alonso Garcia (173291)\\[0.7cm]
    \textbf{Curso:} Desarrollo Movil\\
    \textbf{Docente:} [Completar]\\
    \textbf{Fecha:} \today
    \end{flushleft}

    \vfill
\end{titlepage}

\section{Resumen}
Este informe presenta el desarrollo de una aplicacion movil Android para control academico con tres roles operativos: Administrador, Profesor y Alumno. La solucion integra autenticacion por correo y contrasena en Firebase Authentication, persistencia transaccional en Cloud Firestore y un flujo de asistencia basado en codigos QR. Durante la fase de robustecimiento se migro la logica principal a una arquitectura MVVM, incorporando estados explicitos de interfaz para evitar condiciones de carrera, envios duplicados y visualizacion de listas en estados inconsistentes. El documento incluye diagramas generados directamente en LaTeX y evidencia agrupada por dominio funcional, con analisis tecnico de configuracion, integridad de datos y trazabilidad de operaciones.

\textbf{Palabras clave:} Android, Kotlin, Firebase, Firestore, MVVM, QR, control academico.

\section{Introduccion}
El objetivo del proyecto fue construir una aplicacion movil capaz de administrar materias, asistencia y calificaciones mediante un modelo de roles. La propuesta se implemento con Kotlin y servicios de Firebase para acelerar la integracion de identidad y base de datos. El resultado final prioriza separacion de responsabilidades, trazabilidad de eventos y estabilidad de interfaz en operaciones asincronas.

El repositorio oficial del proyecto es: \url{https://github.com/dontloseyourheadsu/ControlAcademico}.

\section{Contexto del proyecto y organizacion del equipo}
\subsection*{Distribucion de responsabilidades}
\begin{longtable}{>{\raggedright\arraybackslash}p{0.28\linewidth} >{\raggedright\arraybackslash}p{0.22\linewidth} >{\raggedright\arraybackslash}p{0.42\linewidth}}
\toprule
\textbf{Integrante} & \textbf{Rol principal} & \textbf{Responsabilidades} \\
\midrule
Jesus Alvarez Sombrerero & Backend Firebase & Configuracion de autenticacion, colecciones de Firestore, validaciones y soporte de flujo de datos. \\
Camila Garcia Alvarez & Frontend Android & Construccion de pantallas por rol, navegacion, mensajes de estado y consistencia visual. \\
Celeste Isabel Alonso Garcia & QR y persistencia local & Integracion ZXing, gestion de permisos de camara y almacenamiento local con SharedPreferences. \\
\bottomrule
\end{longtable}

\section{Requisitos tecnicos y stack}
\begin{itemize}[leftmargin=*]
    \item Lenguaje principal: Kotlin.
    \item Plataforma objetivo: Android API 24 o superior.
    \item Autenticacion: Firebase Authentication.
    \item Base de datos: Firebase Cloud Firestore.
    \item Persistencia local: SharedPreferences.
    \item Asistencia: generacion y escaneo de QR con ZXing.
\end{itemize}

\section{Arquitectura implementada: MVVM}
La arquitectura implementada es \textbf{MVVM (Model-View-ViewModel)}. Las Activities operan como capa View y se limitan a renderizar estado, reaccionar a eventos de UI y delegar acciones de usuario. Los ViewModels concentran transiciones de estado, validaciones de negocio y coordinacion con el repositorio. El repositorio encapsula la comunicacion con Firestore y desacopla la capa de presentacion del detalle de persistencia.

Esta decision arquitectonica permitio reducir la complejidad ciclomatica de las Activities y mejorar la mantenibilidad del codigo al evitar logica de negocio dispersa. Ademas, la incorporacion de estados \texttt{Loading}, \texttt{Empty}, \texttt{Ready} y \texttt{Submitting} controlo de forma deterministica la habilitacion de controles y la visibilidad de datos.

\section{Diagramas tecnicos (creados en LaTeX)}
\begin{figure}[H]
\centering
\begin{tikzpicture}[
    node distance=2.0cm and 3.0cm,
    box/.style={draw, rounded corners, minimum width=3.2cm, minimum height=1.0cm, align=center, fill=gray!10},
    arr/.style={-{Latex[length=2.5mm]}, thick}
]
\node[box, fill=blue!12] (view) {View\\(Activities)};
\node[box, fill=green!12, right=of view] (vm) {ViewModel\\(estado + eventos)};
\node[box, fill=orange!12, right=of vm] (repo) {Repository\\(FirestoreRepository)};
\node[box, fill=purple!12, below=of repo] (firebase) {Firebase\\Auth + Firestore};
\node[box, fill=yellow!12, below=of view] (model) {Model\\(User, Subject, Grade, Attendance)};

\draw[arr] (view) -- node[midway, above=4pt, fill=white, inner sep=1.5pt]{acciones de usuario} (vm);
\draw[arr] (vm) -- node[midway, above=4pt, fill=white, inner sep=1.5pt]{operaciones de datos} (repo);
\draw[arr] (repo) -- (firebase);
\draw[arr] (repo) |- (model);
\draw[arr] (vm) -- node[midway, left=4pt, fill=white, inner sep=1.5pt]{estado observado} (model);
\draw[arr, dashed] (vm) to[bend left=22] node[midway, above=5pt, fill=white, inner sep=1.5pt]{renderizado por estado} (view);
\end{tikzpicture}
\caption{\textbf{Arquitectura MVVM del sistema}. Flujo principal entre interfaz, estado de presentacion, repositorio y servicios de Firebase.}
\label{fig:mvvm-architecture-latex}
\end{figure}

\begin{figure}[H]
\centering
\begin{tikzpicture}[
    node distance=2.4cm and 2.8cm,
    state/.style={draw, rounded corners, minimum width=3.3cm, minimum height=1.0cm, align=center, fill=gray!10},
    arr/.style={-{Latex[length=2.5mm]}, thick}
]
\node[state, fill=blue!10] (loading) {Loading};
\node[state, fill=green!10, right=of loading] (ready) {Ready};
\node[state, fill=orange!10, below=of ready] (submitting) {Submitting};
\node[state, fill=red!10, below=of loading] (empty) {Empty / Error};

\draw[arr] (loading) -- node[midway, above=4pt, fill=white, inner sep=1.5pt]{consulta exitosa} (ready);
\draw[arr] (loading) -- node[midway, left=4pt, fill=white, inner sep=1.5pt]{sin datos / fallo} (empty);
\draw[arr] (ready) -- node[midway, right=4pt, fill=white, inner sep=1.5pt]{guardar / actualizar} (submitting);
\draw[arr] (submitting) -- node[midway, left=4pt, fill=white, inner sep=1.5pt]{respuesta} (ready);
\draw[arr, dashed] (empty) to[bend right=25] node[midway, below=4pt, fill=white, inner sep=1.5pt]{reintento} (loading);
\end{tikzpicture}
\caption{\textbf{Maquina de estados de UI}. Patron aplicado en pantallas con carga remota para evitar acciones invalidas durante operaciones asincronas.}
\label{fig:ui-state-machine-latex}
\end{figure}

\begin{figure}[H]
\centering
\begin{tikzpicture}[
    node distance=1.7cm and 2.2cm,
    box/.style={draw, rounded corners, minimum width=3.6cm, minimum height=0.95cm, align=center, fill=gray!10},
    arr/.style={-{Latex[length=2.5mm]}, thick}
]
\node[box, fill=blue!12] (auth) {1. Login / Registro};
\node[box, fill=blue!12, below=of auth] (profile) {2. Configuracion de perfil};
\node[box, fill=blue!12, below=of profile] (home) {3. Home por rol};
\node[box, fill=green!12, below left=1.5cm and 2.4cm of home] (admin) {4A. Admin: usuarios + materias};
\node[box, fill=orange!12, below=of home] (prof) {4B. Profesor: notas + asistencia};
\node[box, fill=purple!12, below right=1.5cm and 2.4cm of home] (student) {4C. Alumno: QR + consulta};

\draw[arr] (auth) -- (profile);
\draw[arr] (profile) -- (home);
\draw[arr] (home) -- (admin);
\draw[arr] (home) -- (prof);
\draw[arr] (home) -- (student);
\end{tikzpicture}
\caption{\textbf{Flujo funcional general}. Secuencia de uso desde autenticacion hasta operaciones por rol.}
\label{fig:system-flow-latex}
\end{figure}

\section{Modelo de datos de Firestore}
\subsection*{Colecciones utilizadas}
\begin{itemize}[leftmargin=*]
    \item \textbf{users}: identidad de negocio (uid, nombre, matricula, correo, rol).
    \item \textbf{subjects}: estructura academica (nombre, horario, profesorUid, alumnosUids).
    \item \textbf{grades}: evaluacion por alumno y materia.
    \item \textbf{attendance}: bitacora de asistencia con marca temporal.
\end{itemize}

\subsection*{Control de acceso funcional}
\begin{itemize}[leftmargin=*]
    \item Administrador: gestion de usuarios y materias.
    \item Profesor: registro de calificaciones y asistencia en materias asignadas.
    \item Alumno: consulta de horario, calificaciones y emision de QR.
\end{itemize}

\section{Evidencia visual agrupada y analizada}
\textbf{Nota metodologica:} para mantener legibilidad con alto volumen de capturas, la evidencia se agrupo por dominio funcional y cada bloque se interpreta de forma consolidada.

\subsection*{Bloque 1: Configuracion de Firebase y consistencia de datos}
\begin{figure}[H]
\centering
\begin{subfigure}[t]{0.31\linewidth}
\centering
\EvidenceImage{figures/01-firebase-project.png}
\caption{Proyecto Firebase}
\end{subfigure}
\hfill
\begin{subfigure}[t]{0.31\linewidth}
\centering
\EvidenceImage{figures/02-firebase-auth-users.png}
\caption{Authentication}
\end{subfigure}
\hfill
\begin{subfigure}[t]{0.31\linewidth}
\centering
\EvidenceImage{figures/03-firestore-collections.png}
\caption{Colecciones}
\end{subfigure}

\vspace{0.35cm}

\begin{subfigure}[t]{0.31\linewidth}
\centering
\EvidenceImage{figures/04-firestore-subject-example.png}
\caption{Documento \texttt{subjects}}
\end{subfigure}
\hfill
\begin{subfigure}[t]{0.31\linewidth}
\centering
\EvidenceImage{figures/05-firestore-grade-example.png}
\caption{Documento \texttt{grades}}
\end{subfigure}
\hfill
\begin{subfigure}[t]{0.31\linewidth}
\centering
\EvidenceImage{figures/06-firestore-attendance-example.png}
\caption{Documento \texttt{attendance}}
\end{subfigure}
\caption{\textbf{Evidencia consolidada de backend en Firebase}.}
\label{fig:firebase-group}
\end{figure}

En este bloque se observa la trazabilidad completa entre infraestructura y persistencia. La consola del proyecto confirma una configuracion unica para autenticacion y base de datos, y la lista de usuarios evidencia que los registros creados desde la app fueron sincronizados correctamente con Firebase Authentication. La vista de colecciones valida la separacion del dominio en \texttt{users}, \texttt{subjects}, \texttt{grades} y \texttt{attendance}, lo cual evita mezclar responsabilidades de negocio.

Los documentos mostrados de \texttt{subjects}, \texttt{grades} y \texttt{attendance} reflejan coherencia referencial: la materia contiene profesor y alumnos inscritos, la calificacion vincula alumno-materia-profesor, y la asistencia registra el evento con timestamp. En terminos de configuracion, el proyecto opera con esquema desacoplado por coleccion y consultas orientadas a rol (por UID de profesor/alumno), lo que disminuye joins logicos en cliente y simplifica validaciones. Esta estructura permitio consultas directas por rol y redujo ambiguedades durante lectura de datos en el cliente Android.

\subsection*{Bloque 2: Onboarding y enrutamiento por rol}
\begin{figure}[H]
\centering
\begin{subfigure}[t]{0.31\linewidth}
\centering
\EvidenceImage{figures/07-login.png}
\caption{Login}
\end{subfigure}
\hfill
\begin{subfigure}[t]{0.31\linewidth}
\centering
\EvidenceImage{figures/08-register.png}
\caption{Registro}
\end{subfigure}
\hfill
\begin{subfigure}[t]{0.31\linewidth}
\centering
\EvidenceImage{figures/09-profile-setup.png}
\caption{Perfil inicial}
\end{subfigure}

\vspace{0.35cm}

\begin{subfigure}[t]{0.31\linewidth}
\centering
\EvidenceImage{figures/10-home-admin.png}
\caption{Home Admin}
\end{subfigure}
\hfill
\begin{subfigure}[t]{0.31\linewidth}
\centering
\EvidenceImage{figures/11-home-professor.png}
\caption{Home Profesor}
\end{subfigure}
\hfill
\begin{subfigure}[t]{0.31\linewidth}
\centering
\EvidenceImage{figures/12-home-student.png}
\caption{Home Alumno}
\end{subfigure}
\caption{\textbf{Flujo de autenticacion, perfil y resolucion de rol}.}
\label{fig:onboarding-group}
\end{figure}

Las capturas de acceso muestran un onboarding controlado en tres etapas: autenticacion, captura de metadatos academicos y navegacion al panel correspondiente. La pantalla de perfil opera como punto de normalizacion de datos porque asegura que cada usuario tenga nombre y matricula antes de iniciar operaciones academicas.

Las tres variantes de Home demuestran que la aplicacion resolvio el rol desde Firestore y aplico enrutamiento deterministico. Tecnica y funcionalmente, esto confirma que el estado de sesion no depende solo de Firebase Auth, sino tambien del perfil persistido en \texttt{users}, evitando inconsistencias entre identidad autenticada y permisos de interfaz. Este resultado evita accesos cruzados entre funciones de administracion, docencia y alumno, reforzando el principio de privilegio minimo a nivel de interfaz.

\subsection*{Bloque 3: Operacion por rol (Administrador, Profesor, Alumno)}
\begin{figure}[H]
\centering
\begin{subfigure}[t]{0.31\linewidth}
\centering
\EvidenceImage{figures/13-admin-users-list.png}
\caption{Admin: usuarios y roles}
\end{subfigure}
\hfill
\begin{subfigure}[t]{0.31\linewidth}
\centering
\EvidenceImage{figures/14-admin-create-subject.png}
\caption{Admin: materia y asignaciones}
\end{subfigure}
\hfill
\begin{subfigure}[t]{0.31\linewidth}
\centering
\EvidenceImage{figures/15-admin-subject-created-toast.png}
\caption{Admin: confirmacion}
\end{subfigure}

\vspace{0.35cm}

\begin{subfigure}[t]{0.31\linewidth}
\centering
\EvidenceImage{figures/16-professor-subjects.png}
\caption{Profesor: materias}
\end{subfigure}
\hfill
\begin{subfigure}[t]{0.31\linewidth}
\centering
\EvidenceImage{figures/17-professor-grade-form.png}
\caption{Profesor: calificaciones}
\end{subfigure}
\hfill
\begin{subfigure}[t]{0.31\linewidth}
\centering
\EvidenceImage{figures/18-professor-scan-permission.png}
\caption{Profesor: permiso de camara}
\end{subfigure}
\caption{\textbf{Evidencia de operacion por rol: Administrador y preparacion del flujo docente}.}
\label{fig:roles-group-admin-prof-a}
\end{figure}

\begin{figure}[H]
\centering
\begin{subfigure}[t]{0.31\linewidth}
\centering
\EvidenceImage{figures/19-professor-scan-qr.png}
\caption{Profesor: escaneo QR}
\end{subfigure}
\hfill
\begin{subfigure}[t]{0.31\linewidth}
\centering
\EvidenceImage{figures/20-professor-attendance-saved.png}
\caption{Profesor: asistencia registrada}
\end{subfigure}
\hfill
\begin{subfigure}[t]{0.31\linewidth}
\centering
\EvidenceImage{figures/21-student-qr.png}
\caption{Alumno: QR de asistencia}
\end{subfigure}

\vspace{0.30cm}

\begin{subfigure}[t]{0.31\linewidth}
\centering
\EvidenceImage{figures/22-student-grades.png}
\caption{Alumno: calificaciones}
\end{subfigure}
\hfill
\begin{subfigure}[t]{0.31\linewidth}
\centering
\EvidenceImage{figures/23-student-schedule.png}
\caption{Alumno: horario}
\end{subfigure}
\caption{\textbf{Evidencia de operacion por rol: ejecucion docente y consumo de informacion por Alumno}.}
\label{fig:roles-group-prof-student-b}
\end{figure}

En el tramo de Administrador se observa control efectivo del ciclo de configuracion academica: primero se normalizan roles, despues se crea o actualiza materia, y finalmente se registra una retroalimentacion de exito. Este bloque demuestra que la aplicacion establece dependencias correctas entre identidad (rol), estructura academica (materia) y participantes (profesor/alumnos), lo cual evita que el modulo docente opere sobre grupos no configurados.

En el tramo de Profesor, la evidencia confirma que la operacion academica depende de precondiciones de seguridad y contexto (materia seleccionada, alumno inscrito y permiso de camara). El flujo captura tres controles clave: autorizacion de hardware (camara), validacion de QR contra la materia activa y escritura de asistencia con confirmacion visible. En conjunto, estas capturas muestran un cierre transaccional robusto: la accion de campo se materializa en persistencia verificable y retroalimentacion inmediata al docente.

En el tramo de Alumno, la emision de QR, la consulta de calificaciones y la visualizacion de horario muestran una experiencia de autoservicio acotada al propio usuario. El QR funciona como token de presencia vinculado a la materia seleccionada, mientras que las vistas de calificaciones y horario consumen datos ya normalizados por coleccion. Esta separacion funcional, junto con MVVM y estados de UI, evita condiciones ambiguas en pantalla y mantiene consistencia entre lo visualizado y lo almacenado en Firestore.

\section{Analisis de calidad de software}
Con la migracion a MVVM y el uso de estados explicitos, se observaron mejoras concretas en robustez:

\begin{itemize}[leftmargin=*]
    \item Se eliminaron dependencias directas de negocio dentro de Activities, reduciendo acoplamiento.
    \item Se centralizaron validaciones y transiciones de estado en ViewModels.
    \item Se controlaron respuestas asincronas tardias para evitar sobreescritura de UI con datos obsoletos.
    \item Se habilitaron controles de manera condicional segun estado de carga, evitando acciones invalidas.
\end{itemize}

\section{Conclusiones}
El proyecto alcanza los objetivos funcionales de control academico y agrega un nivel de calidad superior al incorporar arquitectura MVVM, flujo por estados de interfaz y evidencia trazable de cada modulo. La organizacion por roles, junto con persistencia en Firebase y asistencia por QR, ofrece una solucion coherente para escenarios academicos reales. Desde una perspectiva de ingenieria, la separacion View-ViewModel-Repository incrementa mantenibilidad, claridad de responsabilidades y capacidad de evolucion del sistema.

\section{Repositorio y referencias}
\textbf{Repositorio del proyecto:} \url{https://github.com/dontloseyourheadsu/ControlAcademico}

Firebase. (2026). \textit{Firebase documentation}. \url{https://firebase.google.com/docs}

Google. (2026). \textit{Cloud Firestore documentation}. \url{https://firebase.google.com/docs/firestore}

JourneyApps. (2026). \textit{ZXing Android Embedded}. \url{https://github.com/journeyapps/zxing-android-embedded}

Android Developers. (2026). \textit{Guide to app architecture}. \url{https://developer.android.com/topic/architecture}

\end{document}
